\subsection{Summary}

This work set out four design questions. The sensing arrangement of
Section~\ref{subsec:hardware} captures postural, ocular, environmental and temporal indicators
concurrently on a single non-contact node, with two independent channels for viewing distance
(RQ1). Equations~\ref{eq:flags}--\ref{eq:severity} separate transient movement from sustained
exposure through a per-flag hold time, and raise the severity tier when conditions co-occur (RQ2).
The governor of Section~\ref{subsec:governor} converts qualified states into a bounded number of
interventions through deduplication, backoff, escalation and user-controlled suppression (RQ3). The
language model operates inside those bounds as an explanation layer, invoked only when the
qualified flag set changes and constrained to the detected flags (RQ4).

These are design answers, supported by an implementation rather than by measurement. They say
nothing about how accurately the system senses, how well it discriminates ergonomic conditions, or
whether its recommendations change behaviour. All three remain open until the protocol in
Section~\ref{sec:protocol} is executed, and nothing in this paper should be read as evidence for
them.

\subsection{Comparison with Previous Studies}

ERG-AI derives occupational posture from five body-worn accelerometers and generates health
recommendations with a large language model~\cite{ref14}; SitWell classifies nineteen sitting
postures from chair-embedded pressure mats and produces personalised feedback with
GPT-4o~\cite{ref15}. Both demonstrate that sensing and language models can be combined
productively, and both depend on hardware the user must wear or on a specific item of furniture.
Camera-based systems such as PoseX avoid that dependency and reach useful classification accuracy
after a short personal calibration~\cite{ref19}, but remain posture-centric.

The present system differs in two respects. It combines postural, ocular, environmental and
temporal channels on a single non-contact node, and it treats the decision of \emph{whether to
interrupt} as a distinct problem with its own mechanism. In the systems above, an assessment leads
directly to feedback; here, an assessment must first persist, then survive deduplication and
backoff, before it is delivered. No superiority claim is made: the comparison is one of scope, and
the accuracy comparisons that would justify a stronger statement have not been performed.

\subsection{Why Multimodal Monitoring Matters}

The same posture carries different implications depending on its context. A user with a raised head
angle at an appropriate viewing distance, sustained briefly, is not in the same ergonomic situation
as a user who simultaneously holds that angle, sits closer than $30\,\text{cm}$, works under
inadequate lighting, and remains so for twenty minutes. A single-threshold system reports the same
alert in both cases.

Equation~\ref{eq:severity} operationalises this: two individually mild conditions occurring
together raise the severity tier, which is a claim about combinations rather than about any single
signal. Whether that distinction matters in practice --- that is, what proportion of interventions
would arise from combinations rather than from any single threshold --- is measurable from session
logs, and is one of the quantities the protocol in Section~\ref{sec:protocol} is designed to
obtain. It has not been measured here.

\subsection{Detection Is Not Intervention}

Continuous monitoring produces far more detections than a person can usefully act on. A system that
surfaces each one converts a health tool into a source of interruption, and the predictable
consequence is that the user disables it --- at which point measurement accuracy no longer matters
at all. This failure mode receives little attention in the ergonomic monitoring literature, which
generally evaluates detection performance and stops there.

The mechanisms in Section~\ref{subsec:governor} take a concrete position on it. Persistence
filtering removes transient movements that are not exposures. Deduplication on the qualified flag
set prevents an unchanged condition from generating repeated explanations. Backoff lengthens the
interval when a reminder goes unheeded rather than repeating at fixed frequency, and escalation
increases salience once, requiring acknowledgement, instead of increasing frequency indefinitely.
Snooze and do-not-disturb give the user a proportionate response short of disabling the system
entirely.

Whether this produces better adherence than a simpler scheme is an empirical question that has not
been tested here.

\subsection{Practical Implications}

The hardware is inexpensive, non-contact, and mounts on an ordinary monitor without modifying the
workstation, which suits student, home-office and shared-desk settings where instrumented chairs or
body-worn sensors are impractical. Because all thresholds, hold times and governor parameters are
read from a single configuration file and are adjustable from the dashboard, the same deployment
can be retuned for a different user or workspace without redeployment --- a practical requirement
for any ergonomic tool intended to be used by more than one person.
